\section{Extend the caller before changing the queue format}
A logger often needs a policy for a full queue. This implementation reports the
failure without discarding older records. A caller can drop the new record, retry
later, or report lost messages. Overwriting the oldest record would be a different
contract: it would need to advance head across whole encoded records and explain
what happens if several must be discarded. Do not silently add that behavior to
this assignment.

\begin{MARSExercise}{Retry a short destination}
Queue a seven-byte payload in the default ring. Attempt a pop with destination
capacity six, record both return values, and compare the descriptor and destination
with their earlier state. Retry with capacity seven. Explain why the successful
second call returns the same oldest record rather than a later record or a partial
suffix. The supplied demonstration includes this sequence for direct inspection.
\end{MARSExercise}

\begin{MARSExercise}{Stress a non-power-of-two ring}
Use capacity thirteen and start an empty queue at head and tail twelve. Choose
payload lengths whose encoded sizes first leave one free byte, then fill that
byte with an empty record. Pop both records into separate destinations. Trace the
physical prefix and payload slots as well as head, tail, and used. Every state must
satisfy the descriptor invariant, including the two states with equal cursors.
\end{MARSExercise}

\begin{MARSExercise}{Compare contiguous spans}
Replace per-byte wrap tests with two bounded copy spans: the suffix from the current
index to the physical end and, when needed, a prefix beginning at index zero.
Derive each span length before coding. Keep prefix placement separate from payload
placement, preserve zero-length behavior, and compare native counts on wrapped and
unwrapped records. Rerun the complete behavioral checks; a lower count does not
justify writing an adjacent occupied slot.
\end{MARSExercise}

A queue persisted in a file would need validation for damaged lengths and metadata,
a defined byte order, and a recovery policy. A queue used by concurrent threads
would need synchronization and a publication protocol. Those are separate designs.
The present exercise establishes a bounded sequential FIFO whose state and byte
copies can be checked at each call boundary.
